\section{Balanced persistence: affine amplification and explicit counterexample gates}
\label{sec:balanced-persistence}

This section records a further simultaneous attack on the positive and
negative directions.  A route is kept active whenever its missing statement
is merely difficult or unknown.  We mark a construction mechanism as closed
only when a theorem or an explicit counterexample contradicts the statement
that mechanism actually needs.  A local failure is never promoted to a
failure of the surrounding route.

The standard conjecture remains the target throughout.  No existence premise
for a squarefull Pell subsequence, a squarefull Hall remainder, or a large
low-radical affine subfibre is assumed.

\subsection{A complete two-parameter primitive shear}

Fix a primitive positive seed
\[
             a+b=c,\qquad \gcd(a,b)=1,
\]
and put
\[
 P=abc,\qquad R=\rad(P),\qquad
 \rho=\frac{\log P}{\log c},\qquad
 \sigma=\frac{\log R}{\log c}.
\]
After interchanging $a$ and $b$, assume $b\geq c/2$.  The elementary bounds
\begin{equation}\label{eq:bp-P-range}
                  c(c-1)\leq P\leq \frac{c^3}{4}
\end{equation}
will be used repeatedly.

For integers $h,k\geq1$ define
\begin{equation}\label{eq:bp-shear}
 \begin{split}
 U&=1+Ph,\\
 V&=1+P(h+ck),\\
 W&=1+P(h+bk).
 \end{split}
\end{equation}

\begin{theorem}[Primitive affine shear]
\label{thm:bp-affine-shear}
If $\gcd(U,k)=1$, then
\[
                         (aU,bV,cW)
\]
is a positive primitive $abc$ point.  The ordered output determines the
parameter pair $(h,k)$.
\end{theorem}

\begin{proof}
The three identities
\[
 V-U=cPk,\qquad W-U=bPk,\qquad V-W=aPk
\]
and $a+b=c$ give
\[
 aU+bV=cW.
\]
All three cofactors are congruent to $1$ modulo $P$, so they avoid every
prime of the seed.  A common divisor of $U$ and $V$ divides $cPk$ and is
coprime to $cP$; it must therefore divide $k$, contrary to
$\gcd(U,k)=1$.  The same argument gives $\gcd(U,W)=1$.

A common divisor of $V$ and $W$ divides $aPk$.  It avoids $aP$, while
$V\equiv U\pmod{k}$ and hence is coprime to $k$.  Thus $U,V,W$ are pairwise
coprime.  The seed coordinates are pairwise coprime, and all cross terms are
coprime because the cofactors avoid $P$.  This proves primitivity.
Finally $U=(aU)/a$ recovers $h$, and $V-U=cPk$ then recovers $k$.
\end{proof}

The construction is not a thin parametrization.

\begin{lemma}[Uniform admissible count]
\label{lem:bp-admissible-count}
For all sufficiently large $M$,
\[
 \#\{1\leq h,k\leq M:\gcd(1+Ph,k)=1\}\geq \frac{M^2}{8}.
\]
Inside the upper square $\lceil M/2\rceil\leq h,k\leq M$, the count is at
least $M^2/32$.
\end{lemma}

\begin{proof}
A bad pair has a prime $\ell\leq M$ dividing both $k$ and $1+Ph$; such a
prime does not divide $P$.  For each $\ell$, there are at most $M/\ell$
choices of $k$ and at most $M/\ell+1$ choices of $h$.  Hence
\[
 \#\{\hbox{bad pairs}\}
 \leq M^2\sum_{\substack{\ell\leq M\\ \ell\text{ prime}}}\ell^{-2}
       +M\sum_{2\leq n\leq M}n^{-1}
 \leq \frac{3M^2}{4}+M\log M.
\]
Here $\sum_{n\geq2}n^{-2}\leq3/4$.  This proves the first assertion for
large $M$.

For the upper square let $L$ denote its side length.  The contribution of
$\ell$ is at most $(L/\ell+1)^2$, so the total number of bad pairs is at
most
\[
                       \frac{3L^2}{4}+2L\log M+\pi(M).
\]
For large $M$ the final two terms are at most $L^2/8$.  Since $L\geq M/2$,
at least $L^2/8\geq M^2/32$ pairs remain.
\end{proof}

For fixed $K>5$, set
\[
                  M_K=\left\lfloor\frac{c^{K-2}}{4P}\right\rfloor.
\]
Equation \eqref{eq:bp-P-range}, Theorem~\ref{thm:bp-affine-shear}, and
Lemma~\ref{lem:bp-admissible-count} give
\begin{equation}\label{eq:bp-raw-fibre}
 \#\{\hbox{distinct outputs of height at most }c^K\}
              \gg M_K^2\gg c^{2K-10}.
\end{equation}
Indeed,
\[
 cW\leq c+cPM_K(1+b)<c^K.
\]
At $K=8$ this produces at least $c^6/32$ outputs after increasing the fixed
lower threshold for $c$.

\subsection{The actual-radical gate}

The raw count does not inherit the seed's radical saving automatically.  For
$n\geq1$ put
\[
                         E(n)=\frac{n}{\rad(n)}.
\]
Because the cofactors in \eqref{eq:bp-shear} are pairwise coprime and avoid
$P$, radical multiplicativity gives the exact identity
\begin{equation}\label{eq:bp-excess-identity}
 \rad(aUbVcW)=R\rad(UVW)
             =\frac{RUVW}{E(U)E(V)E(W)}.
\end{equation}
Writing $H=cW$, the condition
\[
                       \rad(aUbVcW)<H^\mu
\]
is therefore equivalent to
\begin{equation}\label{eq:bp-excess-gate}
 E(U)E(V)E(W)>\frac{R}{c}\,UVH^{1-\mu}.
\end{equation}
This equivalence is the positive route's exact remaining arithmetic object.

At $K=8$, $\mu=3/4$, and on the upper parameter square, every exceptional
output has
\begin{equation}\label{eq:bp-excess-c14}
                       E(U)E(V)E(W)>\frac{R}{8192}c^{14}.
\end{equation}
To see the direction of the height inequality explicitly, the lower bounds
in the upper square give $UVW\geq c^{20}/8192$, while $H<c^8$, and
\[
 \frac{R}{c}UVH^{1/4}=\frac{RUVW}{H^{3/4}}.
\]

Define the square-root part
\[
 Q(n)=\prod_p p^{\lfloor v_p(n)/2\rfloor}.
\]
Primewise, $Q(n)^2\mid n$ and $Q(n)^2\geq E(n)$.  Thus
\eqref{eq:bp-excess-c14} places every desired output in the tail
\begin{equation}\label{eq:bp-large-square-tail}
                        Q(UVW)\gg R^{1/2}c^7,
\end{equation}
whereas the parameter box has side $M\ll c^4$.  This explains why an
ordinary small-modulus squarefree sieve does not settle either direction.

There is nevertheless a uniform unconditional upper budget.  Let
$\mathcal E_c$ be the $K=8$, $\mu=3/4$ exceptional subfibre.  Then
\begin{theorem}[Three-way exceptional budget]
\label{thm:bp-three-way-budget}
For every fixed subcritical source locus and every $\delta>0$,
\begin{equation}\label{eq:bp-three-way-budget}
 |\mathcal E_c|\ll_\delta
 c^{\min\{12-2\rho,\,6-\sigma,\,8-\rho-\sigma/3\}+\delta},
\end{equation}
uniformly in the seed.
\end{theorem}

\begin{proof}
The first term is the size of the full two-parameter fibre.  The second is
the actual-support $S$-unit estimate from Section~\ref{sec:actual-radical},
applied with target radical below $c^6$ and inherited radical $R$.

For the third term, \eqref{eq:bp-excess-identity} implies
\[
 D:=\rad(UVW)<\frac{c^6}{R}.
\]
Since the three cofactors are pairwise coprime, one of their radicals is
below
\[
                         Z=(c^6/R)^{1/3}=c^{2-\sigma/3}.
\]
The de Bruijn radical count used by
Bernert--Browning--Lichtman--Ter\"av\"ainen~\cite{BBLT} bounds the number of
possible values of that cofactor by
$O_\delta(c^{2-\sigma/3+\delta})$.  The exponent is uniform: for a fixed
source slope $\sigma\in[0,\lambda]$, the relevant de Bruijn exponent ranges
over a compact interval, so a finite mesh and monotonicity reduce the claim
to finitely many fixed estimates.

For fixed $U$, the parameter $h$ is unique and there are at most $M$ choices
of $k$.  Fixed $V$ or $W$ has no larger representation multiplicity.
Since $M\ll c^{6-\rho}$, the third exponent follows.
\end{proof}

\begin{theorem}[Seed-sensitive affine upper gate]
Fix a positive primitive seed $a+b=c$, and put $R=\rad(abc)$.  Let
$X\geq2$ and $0<\mu<1$.  For an arbitrary subset of the admissible positive
integral parameters, let $\mathcal E(X)$ be its set of distinct affine-shear
outputs satisfying $H=cW\leq X$ and $\rad(aUbVcW)<H^\mu$.  Then, for every
$\varepsilon>0$,
\begin{equation}\label{eq:bp-sharp-affine-upper}
 \#\mathcal E(X)\ll_{\mu,\varepsilon}
        R^{-2/3}X^{2\mu/3+\varepsilon}.
\end{equation}
The implied constant is independent of the seed and of the chosen parameter
subset.
\end{theorem}

\begin{proof}
Positivity in $aU+bV=cW=H$ gives $U,V,W\leq H\leq X$.  Exact support
avoidance and cofactor coprimality give
\[
 \rad(aUbVcW)=R\rad(U)\rad(V)\rad(W),
\]
so every exceptional output satisfies
\[
 \rad(U)\rad(V)\rad(W)<X^\mu/R.
\]
All three pair projections $(U,V)$, $(U,W)$, and $(V,W)$ are injective.  The
two-coordinate consequence of the de Bruijn radical count is
\[
 \#\{(m,n):m,n\leq X,\ \rad(m)\rad(n)<Y\}
 \ll_\varepsilon YX^\varepsilon.
\]
Writing $r_1,r_2,r_3$ for the three cofactor radicals, the identity
$(r_1r_2)(r_1r_3)(r_2r_3)=(r_1r_2r_3)^2$ shows that one pair product is below
$(X^\mu/R)^{2/3}$.  Apply the pair count through the corresponding injective
projection and sum over the three choices.  This proves
\eqref{eq:bp-sharp-affine-upper}.
\end{proof}

The relaxed region $\rho\leq3$, $\sigma<1$ leaves every exponent in
\eqref{eq:bp-three-way-budget} above $4$; the infimum of the third term is
$14/3$.  This is a lower envelope of the available upper-budget exponents,
not a pointwise upper bound $|\mathcal E_c|\ll c^{14/3}$.  For a fixed
source slope $\lambda<1$ its corresponding relaxed envelope is
$5-\lambda/3$.

At $\mu=3/4$, Proposition~1.1 of
Bernert--Browning--Lichtman--Ter\"av\"ainen gives exponent
$2\mu/3=1/2$ in target height~\cite{BBLT}.  Their Theorem~1.3 gives the
different uniform exponent $3/5$ for fixed $\mu<1$; at $\mu=3/4$ it is the
weaker of the two applicable estimates.  Hence the following statement would
close the positive route through the subcritical-locus equivalence.

\begin{corollary}[Exact conditional closing step]
\label{cor:bp-affine-closing-step}
Fix $\lambda<1$.  Suppose every sufficiently large seed satisfying
$R<c^\lambda$ has at least $c^{17/4}$ members of $\mathcal E_c$.  Then the
height of that source locus is bounded.
\end{corollary}

\begin{proof}
Every member is a distinct primitive point of height below $c^8$.  For every
$\delta>0$, the cited exceptional-set theorem gives
\[
                         |\mathcal E_c|\ll_\delta c^{4+8\delta}.
\]
Take $\delta=1/64$.  Then $4+8\delta=33/8<17/4$, so comparison with the
hypothesized lower bound bounds $c$.
\end{proof}

The fixed exponent $17/4$ is sufficient but no longer optimal.  By
\eqref{eq:bp-sharp-affine-upper}, the sharp conditional comparison at
$X=c^8$ is
\[
              |\mathcal E_c|\geq R^{-2/3}c^{4+\eta}
              \qquad(\eta>0).
\]
For a seed with $R=c^{\sigma+o(1)}$, the required exponent is therefore
$4-2\sigma/3+\eta$.  This is a conditional lower target, not a consequence of
the upper bound.

There is also a precise local warning.  Conditional on local admissibility at
$p\nmid P$, each of $p^2\mid U,V,W$ has probability $1/(p(p+1))$, but the
three events are pairwise disjoint.  Thus the model asserting independence of
the three same-prime events is refuted.  Finite local data at distinct primes
remain exactly independent by the Chinese remainder theorem.  This local
counterexample closes only that probabilistic model; it does not refute the
affine lower-bound route.

No theorem presently supplies the lower bound in
Corollary~\ref{cor:bp-affine-closing-step}.  The upper budget does not refute
it.  The affine route therefore remains active.

\subsection{Square-root descent on the balancing Pell orbit}

Let
\[
 u_0=0,\quad u_1=1,\quad u_{n+2}=6u_{n+1}-u_n,
\]
and define $x_n$ by
\[
                         x_n^2-8u_n^2=1.
\]
Equivalently,
\[
                         (3+\sqrt8)^n=x_n+u_n\sqrt8.
\]
Introduce the companion coordinates
\[
                         (1+\sqrt2)^n=A_n+B_n\sqrt2.
\]

\begin{theorem}[Coprime square-root factorization]
\label{thm:bp-square-root-factorization}
For every $n\geq1$,
\begin{equation}\label{eq:bp-square-root-factorization}
 x_n=A_n^2+2B_n^2,\qquad u_n=A_nB_n,\qquad
 A_n^2-2B_n^2=(-1)^n,\qquad \gcd(A_n,B_n)=1.
\end{equation}
Consequently $u_n$ is squarefull if and only if both $A_n$ and $B_n$ are
squarefull.
\end{theorem}

\begin{proof}
Squaring $(1+\sqrt2)^n$ and using
$(1+\sqrt2)^2=3+\sqrt8$ gives the first two identities.  Taking norms gives
the third.  A common divisor of $A_n$ and $B_n$ divides the norm, hence is
one.  The prime supports of the two factors are therefore disjoint, which
proves the squarefull equivalence.
\end{proof}

Adding and subtracting the norm identity gives
\begin{equation}\label{eq:bp-adjacent-shapes}
 \left(\frac{x_n-1}{2},\frac{x_n+1}{2}\right)=
 \begin{cases}
  (2B_n^2,A_n^2),&n\text{ even},\\
  (A_n^2,2B_n^2),&n\text{ odd}.
 \end{cases}
\end{equation}
Thus the associated primitive point $Q_n$ is obtained by inserting the unit
between the two entries in \eqref{eq:bp-adjacent-shapes}.

\begin{theorem}[One-half radical slope]
\label{thm:bp-half-slope}
If $u_n$ is squarefull and $c_n$ is the larger nonunit entry of $Q_n$, then
\begin{equation}\label{eq:bp-half-slope}
              \log\rad(Q_n)<\frac12\log c_n+\log2.
\end{equation}
Hence an unbounded squarefull subsequence of $(u_n)$ would disprove the
standard $abc$ conjecture.
\end{theorem}

\begin{proof}
Theorem~\ref{thm:bp-square-root-factorization} makes both factors squarefull,
so
\[
 \rad(Q_n)=\rad(2A_n^2B_n^2)
       \leq2\rad(A_n)\rad(B_n)\leq2\sqrt{A_nB_n}.
\]
For even $n$, $c_n=A_n^2$ and $B_n<A_n$; for odd $n$,
$c_n=2B_n^2$ and $A_n<2B_n$.  Thus $A_nB_n<c_n$ in both cases, proving
\eqref{eq:bp-half-slope}.  Applying standard $abc$ with, for example,
$\eps=1/2$ would bound the heights, contradicting an unbounded subsequence.
\end{proof}

On an even squarefull index, both nonunit entries are $4$-full, so the fixed
signature is $(3,4,4)$ with reciprocal sum $5/6$.  This is an additional
conversion of the same unresolved squarefull premise, rather than a new
source of indices.

\subsection{Valuation saturation and prime-index descent}

For a prime $p$, let $z(p)$ be its rank of apparition in $(u_n)$.  Strong
divisibility gives $p\mid u_n$ if and only if $z(p)\mid n$.  For odd $p$,
the order calculation in the quadratic residue algebra gives
\begin{equation}\label{eq:bp-rank-bound}
                       z(p)\mid p-\left(\frac{2}{p}\right).
\end{equation}
Sanna's exact valuation theorem~\cite{Sanna2016} specializes to
\begin{equation}\label{eq:bp-Sanna}
 v_p(u_n)=
 \begin{cases}
  v_p(u_{z(p)})+v_p(n),&z(p)\mid n,\\
  0,&z(p)\nmid n,
 \end{cases}
 \qquad(p\text{ odd}),
\end{equation}
and
\[
 v_2(u_n)=0\quad(n\text{ odd}),\qquad
 v_2(u_n)=v_2(n)\quad(n\text{ even}).
\]

Equations \eqref{eq:bp-rank-bound}--\eqref{eq:bp-Sanna} yield an exact
saturation rule.  If $p\parallel u_{z(p)}$, $z(p)\mid N$, and $u_N$ is
squarefull, then $p\mid N/z(p)$.  Iterating only the factorizations
\[
 \begin{array}{c|l}
 2&2\cdot3\\
 3&5\cdot7\\
 4&2^2\cdot3\cdot17\\
 5&29\cdot41\\
 6&2\cdot3^2\cdot5\cdot7\cdot11\\
 7&13^2\cdot239
 \end{array}
\]
proves the finite obstruction
\begin{equation}\label{eq:bp-index-saturation}
 u_N\text{ squarefull and }2\mid N
 \quad\Longrightarrow\quad
       22{,}318{,}790{,}340\mid N.
\end{equation}
The factor $13^2$ at rank seven supplies no outgoing exponent-one edge.  This
is a real obstruction to the naive use of primitive-divisor theorems:
\begin{equation}\label{eq:bp-u7}
                         u_7=13^2\cdot239.
\end{equation}
Carmichael's theorem in Yabuta's formulation~\cite{Yabuta2001} provides a
primitive divisor at every relevant index, but \eqref{eq:bp-u7} proves that
``primitive'' does not mean ``valuation one.''

The strongest unconditional descent obtained here reduces the full
squarefull question to prime indices.

\begin{theorem}[Largest-prime descent]
\label{thm:bp-largest-prime-descent}
Suppose $N>1$ and $u_N$ is squarefull.  If $\ell$ is the largest prime
divisor of $N$, then $\ell$ is odd and $u_\ell$ is squarefull.
\end{theorem}

\begin{proof}
If $N$ were a power of two, then $z(3)=2$ and
\eqref{eq:bp-Sanna} would give $v_3(u_N)=1$, a contradiction.  Thus $\ell$
is odd.

Let $q\mid u_\ell$.  Since $u_\ell$ is odd and $\ell$ is prime,
$z(q)=\ell$.  Equation \eqref{eq:bp-rank-bound} gives $\ell\mid q\pm1$.
An odd prime with this property is larger than $\ell$: the first positive
representatives $\ell-1$ and $\ell+1$ are even.  Hence $q\nmid N$ by the
maximality of $\ell$.  Since $\ell\mid N$, valuation lifting gives
\[
                         v_q(u_N)=v_q(u_\ell)+v_q(N)=v_q(u_\ell).
\]
The left side is at least two.  This holds for every prime divisor of
$u_\ell$, proving squarefullness.
\end{proof}

An odd prime $p$ for which $p^2\mid u_{z(p)}$ is a balancing-Wieferich prime.
Write $(1+\sqrt2)^n=A_n+B_n\sqrt2$.  At an odd prime index $\ell$,
$u_\ell=A_\ell B_\ell$ and $\gcd(A_\ell,B_\ell)=1$.  The perfect-power
classifications of Cohn and Ljunggren~\cite{Cohn1996Pell,Cohn2003Pell,Ljunggren1942}
sharpen the prime-index alternative as follows.

\begin{theorem}[Prime-index four-prime packet]
\label{thm:bp-four-prime-packet}
For every odd prime $\ell$, either some prime satisfies
$q\parallel u_\ell$, or $\ell\ne7$ and $u_\ell$ is squarefull with at least
four distinct support primes of rank exactly $\ell$: at least two divide
$A_\ell$ and at least two divide $B_\ell$.  In this second alternative every
such prime has first valuation at least two, and each channel contains a prime
whose first valuation is at least three.
\end{theorem}

\begin{proof}
If no support valuation equals one, the term is squarefull.  Coprimality makes
both channels squarefull.  Cohn's associated-Pell theorem says that
$A_\ell>1$ is not a perfect power.  The standard Pell classifications say the
same for $B_\ell$, except at $\ell=7$; there
$u_7=13^2\cdot239$ and $239\parallel u_7$.  A squarefull integer which is not
a perfect power has at least two support primes and an odd valuation at least
three.  Applying this separately to the two channels proves the packet.
The rank and valuation claims follow from
\eqref{eq:bp-rank-bound}--\eqref{eq:bp-Sanna}.
\end{proof}

The channel orders also force
\[
 q\mid A_\ell\Longrightarrow q\equiv1\text{ or }2\ell+1\pmod{4\ell},
 \qquad
 q\mid B_\ell\Longrightarrow q\equiv1\text{ or }2\ell-1\pmod{4\ell}.
\]
Combining Theorem~\ref{thm:bp-four-prime-packet} with the largest-prime
descent shows that every hypothetical nonunit squarefull balancing term forces
such a packet at an odd prime rank.  The adjacent Pell point consequently has
the sharpened bound
\begin{equation}\label{eq:bp-channel-radical}
 \rad(2A_\ell^2B_\ell^2)
 \leq 2\sqrt{\frac{u_\ell}{4\ell^2-1}}.
\end{equation}
No theorem currently rules out all such packets or constructs infinitely many
of them.  The Pell route remains active.

\subsection{Elliptic and Hall alternatives}

The Mordell curve
\[
                         E:y^2=x^3-2,\qquad P=(3,5)
\]
supplies an independent critical family.  Write
\[
 nP=\left(\frac{\mathsf X_n}{\mathsf D_n^2},
           \frac{\mathsf Y_n}{\mathsf D_n^3}\right)
\]
in lowest terms.  Then
\begin{equation}\label{eq:bp-EDS}
                  \mathsf Y_n^2+2\mathsf D_n^6=\mathsf X_n^3.
\end{equation}
The three entries are pairwise coprime.  Since
\[
 2P=\left(\frac{129}{10^2},\frac{-383}{10^3}\right),
\]
the denominator elliptic divisibility sequence satisfies
$\mathsf D_2\mid\mathsf D_{2k}$, so $\mathsf D_n$ is even at every even
index~\cite{EverestMcLarenWard2006}.  Equation
\eqref{eq:bp-EDS} then has critical signature $(2,6,3)$.  It becomes strict
if $\mathsf Y_n$ or $\mathsf X_n$ is squarefull on an unbounded subsequence,
or if the odd part of $\mathsf D_n$ is squarefull.  Elliptic
primitive-divisor theorems control new support, but not its first valuation,
so none of these upgrade gates is currently refuted.

There is also a size-sensitive Hall gate.

\begin{proposition}[Squarefull Hall gate]
\label{prop:bp-Hall-gate}
Suppose positive primitive data satisfy
\[
                         X^3+K=Y^2,\qquad K^2\leq X,
\]
and $K$ is squarefull.  Then
\begin{equation}\label{eq:bp-Hall-powers}
       \rad(X^3KY^2)^{12}<(Y^2)^{11},
\end{equation}
so an unbounded family of such data would disprove standard $abc$.
\end{proposition}

\begin{proof}
Put $D=\rad(X^3KY^2)$.  Radical submultiplicativity gives
$D\leq X\rad(K)Y$.  Squarefullness and the size condition give
\[
                         \rad(K)^{12}\leq K^6\leq X^3.
\]
Therefore $D^{12}\leq X^{15}Y^{12}$.  Since $K>0$,
$X^3<Y^2$, whence $X^{15}<Y^{10}$.  This proves
\eqref{eq:bp-Hall-powers}.  Its logarithmic slope is strictly below $11/12$.
\end{proof}

Danilov's identity, as recorded by Dujella~\cite{Danilov1982,Dujella2011}, is
\begin{equation}\label{eq:bp-Danilov}
 (z^2+6z+4)^3-(z^2+1)(z^2+9z+19)^2=-27(2z+11).
\end{equation}
Choose positive solutions
\[
 z^2-5w^2=-1,\qquad z\equiv57\pmod{125}.
\]
There are infinitely many: one is $(682,305)$, and multiplication by a fixed
positive power of the norm-one unit $9+4\sqrt5$ preserves the congruence and
gives an unbounded orbit.  Put
\[
 \begin{split}
 A&=z^2+6z+4,\qquad B=z^2+9z+19,\qquad L=2z+11,\\
 X&=A/5,\qquad Y=wB/5,\qquad K=27L/125.
 \end{split}
\]
The congruence gives $125\mid L$, $A\equiv20\pmod{25}$, and $5\mid w$.
Equation \eqref{eq:bp-Danilov} becomes
\begin{equation}\label{eq:bp-Danilov-normalized}
                         X^3+K=Y^2.
\end{equation}
Moreover $\gcd(A,B)=1$ and $\gcd(A,wB)=5$, so $\gcd(X,Y)=1$.  Thus this is
an unbounded primitive Hall family with
\[
       X\sim\frac{z^2}{5},\qquad
       K\sim\frac{54z}{125}
        \sim\frac{54\sqrt5}{125}\sqrt X,\qquad
       \frac{54\sqrt5}{125}<1.
\]
Consequently $K^2<X$ on a sufficiently large tail of the orbit.
The first point is the exact identity
\[
                     93844^3+297=28748141^2.
\]
Here $297=3^3\cdot11$ is not squarefull.  This refutes only the universal
claim that every remainder in the congruence orbit is squarefull; it does not
refute an unbounded squarefull subsequence.  Such a subsequence would satisfy
Proposition~\ref{prop:bp-Hall-gate} after discarding its finite initial segment
and would disprove $abc$.

For each fixed $m\geq2$, forcing the moving factor
$(2z+11)/125$ to be one exact $m$th power leads to the hyperelliptic equation
\[
 (2w)^2=3125t^{2m}-550t^m+25.
\]
The polynomial on the right is squarefree of degree $2m$, so its smooth
completion has genus $m-1\geq1$.  Siegel's theorem yields only finitely many
integral points for that fixed $m$.  This closes the fixed exact-power
mechanism.  It does not close the squarefull route because the square and
cube kernels may both move.

\paragraph{Certified finite computation.}
A reproducible exact audit in the companion computation archive checked both
open squarefull routes without extrapolating from finite data.  For the
balancing recurrence, every one of the $999$ terms $u_n$ with
$2\leq n\leq1000$ has an explicit prime $p$ and residue $q$ such that
\[
                  u_n\equiv pq\pmod {p^2},\qquad 1\leq q<p.
\]
Thus each of those terms is rigorously non-squarefull.  The extended
certified run through $n=2000$ now settles $1997$ of the $1999$ nonunit terms
and leaves precisely
\[
                         \{1873,1951\}
\]
unresolved.  These two indices are neither squarefull hits nor certified
non-hits.  The seven new exclusions use local recursive Pocklington proofs and
two independent recurrences modulo $p^2$.

An exhaustive scan of all $183{,}071$ odd primes $q\leq2{,}500{,}000$ tested
$u_{q-(2/q)}$ modulo $q^3$.  The only balancing-Wieferich hits are
$13,31,1546463$, all of exact depth two; no depth-three hit occurs in this
finite range.  This is pressure against the four-prime packet, not an
infinite nonexistence theorem.

On the Danilov orbit obtained by iterating
$(9+4\sqrt5)^{10}$ from $(682,305)$, the audit verifies the Pell equation,
integrality, primitivity, $X^3+K=Y^2$, and $K^2\leq X$ for
$0\leq t\leq80$; every one of the $81$ remainders also has a certified prime
with $v_p(K)=1$.  The five-small-prime periodic sieve already fails to decide
$t=326$.  A subsequent symbolic prime-square lift eliminates the finite-search
boundary: if $K_t$ is squarefull, then
\begin{equation}\label{eq:bp-Danilov-global-sieve}
 \begin{aligned}
 t&\equiv122136955032565025967809449110840347537827\\
  &\hspace{1em}\pmod{183205432548847538951714173666260521306741}.
 \end{aligned}
\end{equation}
Hence every nonnegative index below the displayed representative is ruled out
as a squarefull index.  The surviving index residue class is nonempty, but the
theorem asserts no squarefull value in it and does not exclude all such values.
The Danilov route remains active.

The Cohn--Nitaj constructions provide genuine unbounded critical $(3,3,3)$
families~\cite{Nitaj1995,Cohn1998}.  Walsh's Theorem~1.1 supplies another
infinite family for an odd prime $p$ under the explicit hypothesis that
\(Y^2=X^3-432p^2\) has positive rank~\cite{Walsh2024}; without a proved rank
statement for the selected $p$, that branch remains conditional.  In each
case the reciprocal sum equals one.  A powerful-coordinate subsequence would
make a strict signature, and no cited theorem either constructs or refutes
that subsequence.  These routes also remain active.

\subsection{Formalization boundary and route ledger}

Seven Lean modules formalize the elementary deterministic core after the
paper proofs were written:
\begin{center}
\begin{tabular}{>{\raggedright\arraybackslash}p{0.36\textwidth}
                >{\raggedright\arraybackslash}p{0.55\textwidth}}
\toprule
Module & Kernel-checked content\\
\midrule
\texttt{PellAdjacentFactor\allowbreak Counterexample20260831}
 & Half-factor existence, primitive adjacent point, radical square bound,
 one-half conductor inequality, fourth-full transfer, and the conditional
 implication to $\neg\texttt{ABCConjecture}$.\\
\texttt{PellSquareRoot\allowbreak Descent20260831}
 & The $(1+\sqrt2)^n$ recurrence, norm, coprimality, exact square shapes,
 fullness splitting, and equivalence with the Pell coordinate.\\
\texttt{PellPrimeIndex\allowbreak Dichotomy20260831}
 & Odd squarefull exponents have depth at least three, the complete
 index-seven obstruction $239\parallel u_7$, and the two-channel numerical
 product bound.\\
\texttt{HallSquarefull\allowbreak Counterexample20260831}
 & The primitive Hall point, the twelfth-power radical inequality, the
 $11/12$ conductor inequality, and the unbounded-family disproof gate.\\
\texttt{AffineShear\allowbreak Amplification20260831}
 & The shear equation, avoidance of the seed support, all cofactor and
 endpoint coprimality statements, the actual \texttt{ABCPoint}, and
 injectivity of the parameter map.\\
\texttt{AffineExcess\allowbreak UpperBound20260831}
 & Injectivity of all three pair projections, same-prime double-divisibility
 exclusion, and the integer geometric-mean lemma.\\
\texttt{DanilovGlobalIndex\allowbreak Sieve20260831}
 & The $11$- and $89$-adic stages, ten prime-square lift packets, exact CRT,
 and the global progression and lower-index exclusion in
 \eqref{eq:bp-Danilov-global-sieve}.\\
\bottomrule
\end{tabular}
\end{center}

The analytic BBLT and $S$-unit estimates; the Cohn--Ljunggren perfect-power
classifications; Sanna's valuation theorem; Carmichael--Yabuta and Siegel;
elliptic-divisibility and primitive-divisor results; the Danilov
parametrization, primitivity, and orbit unboundedness; and the Cohn--Nitaj and
conditional Walsh existence results all remain outside this Lean increment.
They are proved on paper or cited, rather than inserted as project axioms.

The resulting route ledger is deliberately narrow:
\begin{itemize}
\item the affine repeated-prime lower bound, simultaneous squarefull Pell
coordinates, exclusion or construction of the forced four-prime
balancing-Wieferich packets, the surviving Danilov residue class, and powerful
elliptic-coordinate subsequences are \emph{active};
\item a fixed strict generalized-Fermat equation, a finite residual-kernel
packet, a primitive polynomial-full identity, and the fixed exact-power
Danilov specialization are closed only as the stated construction mechanisms;
\item the example $K=297$ refutes ``every Danilov remainder is squarefull''
but leaves all unbounded-subsequence variants open; the same-prime independent
affine model is also refuted, without closing the affine route.
\end{itemize}
No unconditional closed term of type \texttt{ABCConjecture}, and no rigorous
unconditional disproof, is obtained in this increment.
